\subsection{A response in the right direction can still miss the target}
All 1,008 outputs yield valid measurements. With the full-frame
reference target, every model's aligned-response interval lies above zero
(Table~\ref{tab:specificity-models}). A purely common additive response is
therefore insufficient to explain the named conditions under the stated
inference model. Yet positive alignment does not ensure low error. GPT Image 2
nearly matches the reference-aligned strength, $\beta=.999\;[.893,1.104]$,
while its total error is $D=1.226$: it adds substantial differences beyond
that aligned component.

\begin{table}[tbp]
\centering\small\setlength{\tabcolsep}{4pt}
\begin{tabularx}{\linewidth}{@{}>{\raggedright\arraybackslash}p{.24\linewidth}>{\raggedright\arraybackslash}Xr>{\raggedright\arraybackslash}X@{}}
\toprule
Model & Aligned response $\beta$ & \shortstack{Uncalibrated\\error $D$} & $\Delta D$ versus FLUX\\
\midrule
GPT Image 1 & $0.938\;[0.709,1.168]$ & $1.572$ & $0.771\;[-0.037,1.580]$\\[3pt]
GPT Image 2 & $0.999\;[0.893,1.104]$ & $1.226$ & $0.425\;[-0.161,1.011]$\\[3pt]
GPT Image 2.5 Flare & $0.772\;[0.680,0.864]$ & $1.738$ & $0.937\;[0.256,1.618]$\\[3pt]
GPT Image 2.5 Sunburst & $0.824\;[0.680,0.968]$ & $1.641$ & $0.840\;[0.058,1.623]$\\[3pt]
Nano Banana 2 & $0.442\;[0.256,0.629]$ & $1.109$ & $0.309\;[-0.847,1.464]$\\[3pt]
FLUX.2 Max & $0.470\;[0.239,0.701]$ & $0.801$ & $0$ (baseline)\\[3pt]
\bottomrule\end{tabularx}
\caption{Agreement with the full-frame reference target. $\beta=1$ matches its aligned amplitude; larger is not necessarily better. Expected $D$ is 0 for exact agreement and 1 for no painter distinctions. Positive $\Delta D$ favors FLUX.2 Max. Table brackets give approximate 95\% simultaneous intervals from the prespecified family of six aligned responses and 15 scene-paired model contrasts. FLUX's lowest $D$ estimate has an unadjusted 95\% interval $[0.586,1.016]$, spanning the no-distinction value 1; this interval is descriptive.}
\label{tab:specificity-models}\end{table}


FLUX.2 Max has the lowest error estimate, $.801$, despite a smaller aligned
response, $.470$. Adjusted comparisons resolve higher error for Flare and Sunburst
than for FLUX; the other 13 intervals include zero (Appendix~\ref{app:results},
Table~\ref{tab:all-specificity-pairs}).
These use full-frame references;
the Sunburst result changes after source correction
(Section~\ref{sec:source-correction}). No model is established to beat the
no-contrast benchmark of one. Thus the study detects a
reference-aligned response more clearly than it establishes low total mismatch.

In Figure~\ref{fig:specificity-geometry}, C\'ezanne remains visibly separated,
while generated Monet and Sisley means often lie close together. The M/S/P detail
boxes magnify the nearby positions. The axes combine the 31 standardized features, with directions fitted
only to the four reference means. Together they show 86.99\% of variation among
those means; the pair analysis below checks all 31 features.

\begin{figure}[tbp]
\centering\includegraphics[width=\linewidth]{specificity_geometry.pdf}
\caption{Matching the labeled painter pattern. Crosses mark reference contrasts;
open circles average generated contrasts over all 14 scenes and both repeats.
Shorter matching-label lines indicate less projected mismatch;
faint dots show centered scene/repeat contrasts.
The outlined region is enlarged in the adjacent Monet (M), Sisley (S) and
Pissarro (P) detail at a common scale; C\'ezanne
remains in the overview. Each view uses common axes across models. Agreement
in these two projected axes does not establish agreement in all 31 features.}
\label{fig:specificity-geometry}
\end{figure}
\FloatBarrier

\subsection{Most scene-averaged change is common to the painter names}\label{sec:shared-change}
After averaging scenes, named-minus-artist-free shifts comprise a four-name mean
and painter-specific departures. The mean includes the oil-painting clause and
any common naming response: 82.5--95.7\% of repeat-corrected squared change
(Table~\ref{tab:shared-change}). Averaging emphasizes consistent shifts;
splitting within scenes retains scene-specific variation and gives 88.6--97.2\%
(Appendix~\ref{app:estimation}). $\beta$ measures departures along the reference contrasts.

\begin{table}[htbp]\centering\small\setlength{\tabcolsep}{3pt}
\begin{tabularx}{\linewidth}{@{}l*{6}{>{\centering\arraybackslash}X}@{}}
\toprule
& \shortstack{GPT Image\\1} & \shortstack{GPT Image\\2} & Flare & Sunburst & \shortstack{Nano\\Banana 2} & \shortstack{FLUX.2\\Max}\\\midrule
Shared (\%) & 82.5 & 85.8 & 84.8 & 83.9 & 88.1 & 95.7\\
Between-name (\%) & 17.5 & 14.2 & 15.2 & 16.1 & 11.9 & 4.3\\
\midrule
Corrected cosine & 0.689 & 0.774 & 0.766 & 0.837 & 0.930 & 0.916\\
Squared-size ratio & 1.986 & 1.960 & 2.816 & 3.356 & 3.692 & 4.138\\
\bottomrule\end{tabularx}
\caption{Shares of scene-averaged, repeat-corrected squared change from artist-free to named prompts, including the painting clause. Shared is the four-name mean shift; between-name is departure from it. The lower rows give retrospective comparisons of the shared named shift with the generic oil-painting shift, using the same artist-free baseline. The noise-corrected cosine measures their directional agreement (1 means parallel); the ratio divides their repeat-corrected squared magnitudes, shared by generic.}
\label{tab:shared-change}\end{table}


The shared named shift points broadly along the generic oil-painting shift
(noise-corrected cosine $.689$--$.930$; 1 means parallel), but its repeat-corrected
squared magnitude is 1.96--4.14 times as large. Both use scene averages and the
same artist-free baseline; Table~\ref{tab:shared-change} lists each model's values.
Thus, a response shared across painter names need not be identical to the response
to a generic painting instruction. The generic shift and the additional shared
naming shift can reinforce or oppose one another, so their squared magnitudes
do not add as independent contributions (Appendix~\ref{app:estimation}).
\FloatBarrier

\subsection{The aggregate response conceals uneven painter distinctions}
Retrospective pair diagnostics extend the lightness example in
Section~\ref{sec:contrast-intuition} to all 31 features (Figure~\ref{fig:artist-pairs}).
Monet--Sisley aligned responses range from $-.249$ to $.129$ in five models,
versus $.402$ for GPT Image 2. Yet GPT Image 1 has lower estimated error for this pair
($.56$ versus $1.16$): stronger alignment does not ensure lower mismatch.
Omitting C\'ezanne lowers FLUX's overall response from $.470$ to $.096$ and
GPT Image 2's from $.999$ to $.651$.

Swapping Monet and Sisley lowers error for Flare, Sunburst and Nano Banana 2.
Reversing a negatively aligned pair difference lowers error without changing its size.

\begin{figure}[!htbp]
\centering\includegraphics[width=\linewidth]{specificity_artist_pairs.pdf}
\caption{Descriptive painter-pair estimates (no pairwise tests). M, S, P and C denote Monet, Sisley,
Pissarro and C\'ezanne. Pair aligned response 0 means no aligned component,
1 matches reference strength, values above 1 exaggerate it, and negative values
indicate the opposite direction.
Expected error is 0 for exact agreement and 1 for no distinction. Each pair uses
all 31 features and is normalized by its squared reference distance.
Negative error can arise from repeat correction,
not better-than-exact agreement.}
\label{fig:artist-pairs}
\end{figure}
\FloatBarrier

\subsection{Scene averaging and contrast magnitude change the ordering}
We compare fit scene by scene, after averaging scenes, and after held-out
rescaling of painter contrasts (Table~\ref{tab:review-calibration}).

\emph{Averaging scenes.} Nano Banana 2's mean painter contrasts have lower error than FLUX's
($.723$ versus $.761$). Evaluating each scene adds a larger variation component
for Nano Banana 2 ($.387$ versus $.040$), reversing that order.
$D$ penalizes variation around the pooled historical target; averaging can
cancel scene-dependent differences.

\begin{table}[!htbp]\centering\small
\begin{tabularx}{\linewidth}{@{}l*{3}{>{\centering\arraybackslash}X}@{\hspace{2em}}>{\centering\arraybackslash}X@{}}\toprule
Model & \shortstack{Uncalibrated\\scene error\\$D$} & \shortstack{Error after\\averaging scenes\\$D_{\rm aggregate}$} & \shortstack{Scene\\variation\\$V_{\rm scene}$} & \shortstack{Calibrated\\scene error\\$D_{\rm held}$}\\
\cmidrule(r){1-4}\cmidrule(l){5-5}
GPT Image 1 & 1.572 & 1.239 & 0.333 & 0.645\\
GPT Image 2 & 1.226 & 1.044 & 0.182 & \textbf{0.554}\\
2.5 Flare & 1.738 & 1.483 & 0.255 & 0.741\\
2.5 Sunburst & 1.641 & 1.337 & 0.305 & 0.706\\
Nano Banana 2 & 1.109 & \textbf{0.723} & 0.387 & 0.832\\
FLUX.2 Max & \textbf{0.801} & 0.761 & 0.040 & 0.714\\
\bottomrule\end{tabularx}
\caption{Retrospective error decomposition and calibration against full-frame references. $D=D_{\rm aggregate}+V_{\rm scene}$. Calibration rescales centered generated contrasts using the other 13 scenes and evaluates them on the omitted scene; it does not produce new images. Bold marks each error column's minimum, not significance.}
\label{tab:review-calibration}\end{table}


\emph{Rescaling contrasts.} GPT Image 2's aligned response is near one, but its
squared contrast size exceeds twice the reference's ($Q=2.223$, where 1 matches the reference size;
Appendix~\ref{app:magnitude-coverage}).
Shrinking weakens its reference-aligned component, but reduces off-target
differences enough to lower the total error estimate. GPT Image 2's fitted multipliers
scale every contrast coordinate by about 0.45, using only the other 13 scenes each time.

After calibration, GPT Image 2 has the lowest error estimate, $.554$,
compared with FLUX's $.714$; it has lower scores in 12 of 14 held-out folds.
Deleting any scene and refitting the entire procedure preserves the mean
advantage (Appendix~\ref{app:magnitude-coverage}).
The folds overlap, so their win count does not establish statistical superiority.

\subsection{Which comparisons survive reference changes?}\label{sec:source-correction}
For the uncalibrated scene error $D$, FLUX retains the lowest point
estimate across scene deletions, measurement windows, reference targets and
feature-weighting sensitivities (Appendices~\ref{app:declared-sensitivities}
and~\ref{app:diagnostics}).

Source handling has a more consequential effect on statistical comparisons.
AI assistants inspected all reference and development reproductions, selecting
crops that exclude peripheral artifacts in 90 reference and 41 development images.
The audit used no human verification. After cropping and refitting development scaling,
FLUX's error estimate changes from $.801$ to $.872$. Corrected aligned-response
intervals remain above zero for all six models, and GPT Image 2 retains the lowest
calibrated estimate.

Table~\ref{tab:source-comparison} makes the changed comparisons explicit. With
corrected regions and refitted scaling, the Sunburst interval crosses zero,
Flare still has higher error than FLUX and the GPT Image 1 interval moves above zero. The latter
is a retrospective sensitivity, not a new confirmatory finding. Source correction
preserves FLUX's lowest point estimate while changing the evidence for particular
model differences.

\begin{table}[!htbp]\centering\small
\begin{tabularx}{\linewidth}{@{}Xcc@{}}\toprule
Model comparison & \shortstack{Full-frame difference\\{[adjusted interval]}} & \shortstack{Crop + scaling correction\\Difference [adjusted interval]}\\\midrule
GPT Image 1 $-$ FLUX & $0.771\;[-0.037,1.580]$ & $0.758\;[0.026,1.491]$\\[2pt]
2.5 Flare $-$ FLUX & $0.937\;[0.256,1.618]$ & $0.864\;[0.193,1.535]$\\[2pt]
2.5 Sunburst $-$ FLUX & $0.840\;[0.058,1.623]$ & $0.705\;[-0.026,1.436]$\\[2pt]
\bottomrule\end{tabularx}
\caption{Uncalibrated error differences before and after source correction, using audited painting regions and refitted development scaling. Positive differences favor FLUX.2 Max. Full-frame intervals use the 95\% simultaneous procedure across 21 endpoints; corrected intervals reuse that procedure retrospectively and add no confirmatory tests.}
\label{tab:source-comparison}\end{table}


\begin{samepage}
A separate check uses class-specific reference means for 11 water, built and
land briefs, excluding three mixed briefs. Replacing title-derived reference
labels with visual labels retains FLUX's lowest error estimate.
Appendix~\ref{app:diagnostics} gives the controls and audit.\par
\end{samepage}
